% Standalone Helix: 27 Oct 2026 Research Meeting
% Visual style follows the 21 Oct 2026 light Madrid decks.
% Facts are taken from the local ASPLOS '25 PDF only.
% Compile with pdfLaTeX from this directory. Keep Helix/ beside this file.
\documentclass[aspectratio=169,10pt]{beamer}
\usepackage[english]{babel}
\usepackage[utf8]{inputenc}
\usepackage[T1]{fontenc}
\usepackage{lmodern}
\usepackage{amsmath,amssymb}
\usepackage{booktabs,tabularx,array}
\usepackage{graphicx}
\graphicspath{{./}{figures/}}
\usepackage{xcolor}
\usepackage{tikz}
\usetikzlibrary{arrows.meta,positioning,calc,shapes.geometric}
\usepackage{hyperref}

\usetheme{Madrid}
\usecolortheme{default}
\definecolor{Ink}{HTML}{243748}
\definecolor{Muted}{HTML}{667B8B}
\definecolor{Accent}{HTML}{347E74}
\definecolor{AccentBlue}{HTML}{507DAD}
\definecolor{SoftMint}{HTML}{E6F4EE}
\definecolor{SoftBlue}{HTML}{E9F2F9}
\definecolor{SoftLilac}{HTML}{F1EEF8}
\definecolor{Paper}{HTML}{FBFCFE}
\definecolor{Divider}{HTML}{D8E4EB}

\setbeamercolor{background canvas}{bg=Paper}
\setbeamercolor{normal text}{fg=Ink}
\setbeamercolor{structure}{fg=Accent}
\setbeamercolor{frametitle}{fg=Ink,bg=SoftBlue}
\setbeamercolor{title}{fg=Ink}
\setbeamercolor{subtitle}{fg=Muted}
\setbeamercolor{block title}{fg=Ink,bg=SoftBlue}
\setbeamercolor{block body}{fg=Ink,bg=white}
\setbeamercolor{item}{fg=Accent}
\setbeamercolor{caption name}{fg=Muted}
\setbeamercolor{footline}{fg=Muted,bg=Paper}
\setbeamerfont{frametitle}{size=\large,series=\bfseries}
\setbeamerfont{footline}{size=\tiny}
\setbeamerfont{block title}{series=\bfseries}
\setbeamertemplate{headline}{}
\setbeamertemplate{navigation symbols}{}
\setbeamertemplate{blocks}[rounded][shadow=false]
\setbeamertemplate{itemize items}[circle]
\setbeamertemplate{itemize subitem}[circle]
\setbeamertemplate{footline}{%
  \leavevmode\hbox{%
    \begin{beamercolorbox}[wd=\paperwidth,ht=.30cm,dp=.16cm,rightskip=.60cm plus 1fil]{footline}%
      \hfill\insertframenumber\,/\,\inserttotalframenumber
    \end{beamercolorbox}}}
\setbeamertemplate{frametitle}{%
  \nointerlineskip
  \begin{beamercolorbox}[wd=\paperwidth,ht=.72cm,dp=.16cm,leftskip=.55cm,rightskip=.40cm]{frametitle}%
    \usebeamerfont{frametitle}\insertframetitle
  \end{beamercolorbox}}

\newcommand{\smalllabel}[1]{{\scriptsize\bfseries\color{Accent}\MakeUppercase{#1}}}
\newcommand{\hint}[1]{{\scriptsize\color{Muted}#1}}
\newcommand{\takeaway}[1]{%
  \begin{beamercolorbox}[rounded=true,sep=0.12cm,wd=\linewidth]{block title}%
    {\small\bfseries #1}
  \end{beamercolorbox}}
\newcommand{\stepcard}[3]{%
  \begin{beamercolorbox}[rounded=true,sep=0.10cm,wd=\linewidth]{block body}%
    {\scriptsize\bfseries\color{Accent}#1}\par\vspace{.4mm}
    {\footnotesize\bfseries #2}\par\vspace{.4mm}
    {\scriptsize\color{Muted}#3}%
  \end{beamercolorbox}}

\newcommand{\figureplaceholder}[3]{%
  \begingroup\centering
  \IfFileExists{#1}{%
    \includegraphics[width=.98\linewidth,height=#2,keepaspectratio]{#1}%
  }{%
    \IfFileExists{figures/#1}{%
      \includegraphics[width=.98\linewidth,height=#2,keepaspectratio]{figures/#1}%
    }{%
      \begin{tikzpicture}
        \node[draw=AccentBlue!55,dashed,line width=.7pt,rounded corners=4pt,
              fill=SoftBlue!60,minimum height=#2,minimum width=.95\linewidth,
              text width=.85\linewidth,align=center,inner sep=2mm] {
           {\color{AccentBlue}\bfseries FIGURE PLACEHOLDER}\\[2.2mm]
           {\color{Ink}\small #3}\\[1.5mm]
           {\color{Muted}\scriptsize\ttfamily\detokenize{#1}}};
      \end{tikzpicture}%
    }%
  }\par\endgroup}

\newcommand{\src}[1]{\vspace{0.05cm}\hint{Source: #1}}

\title[Helix]{Helix: Serving LLMs over Heterogeneous GPUs and Network via Max-Flow}
\subtitle{Joint Model Placement and Per-Request Scheduling}
\author{Yu-Jui Liang}
\institute{National Taiwan University \quad | \quad Academia Sinica}
\date{27 Oct 2026}
\begin{document}

% ----- SLIDE 01 -----
{\setbeamertemplate{footline}{}
\begin{frame}[plain]
  \vspace*{0.85cm}
  \smalllabel{Research meeting}\par
  \vspace{.42cm}
  {\raggedright\color{Ink}\LARGE\bfseries
    Helix: Serving Large Language Models over\\[0.18cm]
    Heterogeneous GPUs and Network via Max-Flow\par}
  \vspace{.32cm}
  {\color{Muted}\large Joint Model Placement and Per-Request Scheduling\par}
  \vspace{.85cm}
  {\bfseries Yu-Jui Liang}\par\vspace{.06cm}
  {\small\color{Muted}Advisor: Hsiang-Yun Cheng}\par\vspace{.04cm}
  {\small\color{Muted}National Taiwan University \quad | \quad Academia Sinica \quad | \quad 27 Oct 2026\par}
  \vfill
  {\color{Accent}\rule{\textwidth}{.7pt}}\par\vspace{.08cm}
  {\scriptsize\color{Muted}Paper: Mei et al., ASPLOS 2025. Local PDF is the sole factual source.}
\end{frame}}

% ----- SLIDE 02 -----
\begin{frame}[t]{Presentation Roadmap}
  \smalllabel{Helix / Overview}\par
  \vspace{.18cm}
  \begin{enumerate}
    \setlength{\itemsep}{.16cm}
    \item \textbf{Background / Motivation}\par
      {\small\color{Muted}Why mixed GPUs and wide-area links break uniform pipelines.}
    \item \textbf{Methodology}\par
      {\small\color{Muted}Max-flow placement, MILP constraints, and per-request pipelines.}
    \item \textbf{Experimental Results}\par
      {\small\color{Muted}Single-cluster prototype, plus simulated geo-distributed and 42-node runs.}
    \item \textbf{Discussion / Conclusion}\par
      {\small\color{Muted}What the ratios do and do not say, next to ThunderServe, Hetis, and Tessera.}
  \end{enumerate}
  \src{ASPLOS '25 Helix, local PDF. The comparison page is our synthesis.}
\end{frame}

% ----- SLIDE 03 -----
\begin{frame}[t]{Why Heterogeneous and Geo-Distributed Serving?}
  \smalllabel{Helix / Background}\par
  \vspace{.08cm}
  {\footnotesize\textbf{Table 3.} Datasheet properties. SXM for H100 and A100. Not measured serving throughput.}\par
  \vspace{.04cm}
  {\centering\footnotesize
  \begin{tabular}{lcccc}\toprule
    GPU & FP16 (TFLOP/s) & Memory & Bandwidth & Price (USD)\\\midrule
    H100 & 1979 & 80 GB & 3350 GB/s & 25k--40k\\
    A100 & 312 & 40 GB & 1555 GB/s & 10k--15k\\
    L4 & 242 & 24 GB & 300 GB/s & $\sim$3k\\
    T4 & 65 & 16 GB & 300 GB/s & $\sim$1k\\\bottomrule
  \end{tabular}\par}
  \vspace{.08cm}
  \begin{itemize}\footnotesize
    \setlength{\itemsep}{.05cm}
    \item Table 1: if half the memory stores parameters, LLaMA-2 70B needs 12 L4s, 7 A100s, or 4 H100s.
    \item Table 4: in one us-east window the quota was 8 A100 + 8 L4 + 16 T4, but only 16 T4s were deployed.
    \item The paper's claim: eight L4s have FP16 throughput comparable to one H100, with more memory and less power.
  \end{itemize}
  \vspace{.06cm}\takeaway{Scarcity, not a universal speed ranking, is why Helix aggregates mixed GPUs across regions.}
  \src{Sec.~1, Tables 1--4. Eight-L4 comparison is a datasheet argument.}
\end{frame}

% ----- SLIDE 04 -----
\begin{frame}[t]{Parallelism Across Nodes and GPUs}
  \smalllabel{Helix / Background}\par
  \vspace{.12cm}
  \begin{columns}[T,onlytextwidth]
    \begin{column}{.49\textwidth}
      \stepcard{TENSOR PARALLELISM}{Inside one node}{Each layer needs two collectives. Helix does not use TP across the high-latency network.}
      \vspace{.12cm}
      \stepcard{PIPELINE PARALLELISM}{Across nodes}{Only the activation crosses a stage boundary. Equal cuts leave pipeline bubbles on slow GPUs.}
    \end{column}
    \begin{column}{.49\textwidth}
      \stepcard{PETALS}{Decentralized, greedy}{A new machine takes the least-loaded stage. There is no global plan for a known cluster.}
      \vspace{.12cm}
      \stepcard{SWARM}{Even cuts, local routing}{The model is partitioned evenly. The next stage is chosen from real-time throughput. It is a training method, reused here as a baseline.}
    \end{column}
  \end{columns}
  \vspace{.16cm}\takeaway{Helix keeps TP inside a node and uses PP across nodes, then refuses to cut the model into equal stages.}
  \src{Sec.~2.2 and Sec.~3. Swarm is Ryabinin et al.; Petals is Borzunov et al.}
\end{frame}

% ----- SLIDE 05 -----
\begin{frame}[t]{Uniform Model Placement Is Sub-Optimal}
  \smalllabel{Helix / Motivation}\par
  \vspace{.04cm}
  \figureplaceholder{Helix/Helix_Figure_01_Model_Placement_Examples.png}{3.35cm}{Figure 1}
  {\centering\hint{Fig.~1a--c. Illustrative example, not the 24-node cluster. Order: A100 $>$ L4 $>$ T4.}\par}
  \vspace{.02cm}
  {\footnotesize (b) Equal cuts still wait on the slower stage. (c) Balanced FLOPs still cross the slow inter-region link on every request.}
  \vspace{.04cm}\takeaway{A balanced FLOP cut is not a maximum-throughput serving plan.}
  \src{Sec.~3.1 and Sec.~4.2, Fig.~1a--c. PDF p.~4.}
\end{frame}

% ----- SLIDE 06 -----
\begin{frame}[t]{Network-Unaware Routing Creates Congestion}
  \smalllabel{Helix / Motivation}\par
  \vspace{.04cm}
  \figureplaceholder{Helix/Helix_Figure_01_Model_Placement_Examples.png}{3.35cm}{Figure 1}
  {\centering\hint{Fig.~1c--d. Same illustrative cluster as the previous slide.}\par}
  \vspace{.04cm}
  \begin{itemize}\footnotesize
    \setlength{\itemsep}{.05cm}
    \item In (c) the A100's send side congests: every activation still crosses into Region 2.
    \item (d) runs two T4s beside the A100, so less traffic uses the slow link.
    \item Partition, placement, and routing have to be chosen together.
  \end{itemize}
  \vspace{.06cm}\takeaway{Network capacity is part of placement, not a detail added after the GPUs are assigned.}
  \src{Sec.~4.2, Fig.~1c--d. PDF p.~4.}
\end{frame}

% ----- SLIDE 07 -----
\begin{frame}[t]{Helix Design Requirements}
  \smalllabel{Helix / Problem}\par
  \vspace{.14cm}
  \begin{columns}[T,onlytextwidth]
    \begin{column}{.32\textwidth}
      \stepcard{1 / PLACEMENT}{One plan per cluster}{Assign each node a contiguous layer range that respects GPU memory and the links it will use.}
    \end{column}
    \begin{column}{.32\textwidth}
      \stepcard{2 / ESTIMATE}{Tokens per second}{Score a placement by how many tokens can move from the coordinator, through the GPUs, and back.}
    \end{column}
    \begin{column}{.32\textwidth}
      \stepcard{3 / SCHEDULE}{One path per request}{Do not freeze a few disjoint pipelines. Let requests share nodes when the flow says they should.}
    \end{column}
  \end{columns}
  \vspace{.22cm}
  \begin{itemize}\footnotesize
    \setlength{\itemsep}{.06cm}
    \item Offline placement and online scheduling are different timescales. They share one graph.
    \item KV-cache capacity is not a third MILP objective. It is enforced later by masking nodes.
  \end{itemize}
  \vspace{.10cm}\takeaway{Helix separates a slow placement search from a fast per-request routing decision.}
  \src{Sec.~3 and Sec.~5. The KV split is the paper's design, stated across Sec.~4.4 and Sec.~5.2.}
\end{frame}

% ----- SLIDE 08 -----
\begin{frame}[t]{Helix System Overview}
  \smalllabel{Helix / System}\par
  \vspace{.02cm}
  \figureplaceholder{Helix/Helix_Figure_03_System_Overview.png}{3.05cm}{Figure 3}
  {\centering\hint{Fig.~3. Coordinator plans once; workers execute the assigned layers.}\par}
  \vspace{.04cm}
  \begin{columns}[T,onlytextwidth]
    \begin{column}{.49\textwidth}
      {\footnotesize\begin{itemize}\setlength{\itemsep}{.03cm}
        \item Profiler collects cluster info. The planner solves placement once per cluster.
        \item A new request gets a pipeline and is sent to the first node only.
      \end{itemize}}
    \end{column}
    \begin{column}{.49\textwidth}
      {\footnotesize\begin{itemize}\setlength{\itemsep}{.03cm}
        \item Each node runs its layers and forwards the activation.
        \item The last node returns the token. The next token reuses the same pipeline.
      \end{itemize}}
    \end{column}
  \end{columns}
  \vspace{.04cm}\takeaway{The coordinator owns placement and scheduling. Compute nodes do not re-plan the route.}
  \src{Sec.~5, Fig.~3. PDF p.~7.}
\end{frame}

% ----- SLIDE 09 -----
\begin{frame}[t]{From Model Placement to a Flow Network}
  \smalllabel{Helix / Formulation}\par
  \vspace{.10cm}
  \begin{itemize}\footnotesize
    \setlength{\itemsep}{.08cm}
    \item A cluster has one coordinator \(h\) and compute nodes \(C\). A multi-GPU machine is one logical node: compute and VRAM are summed.
    \item Placement \(\Psi\) gives each node a contiguous layer range. Scheduling \(\phi\) gives each request an ordered node sequence.
    \item A common score keeps only the slowest stage, \(\min_i \sum_j \mathrm{capacity}(j)\cdot \mathbf{1}[i\in\Psi(j)]\).
    \item Fig.~1c has balanced stage FLOPs and is still limited by the wide-area link. The min-stage score cannot see that.
  \end{itemize}
  \vspace{.10cm}\takeaway{Helix scores a placement by the max flow of a graph, in tokens per second.}
  \src{Sec.~4.1--4.2. The min-stage formula is the paper's description of prior heuristics.}
\end{frame}

% ----- SLIDE 10 -----
\begin{frame}[t]{Max-Flow Graph Abstraction}
  \smalllabel{Helix / Formulation}\par
  \vspace{.02cm}
  \begin{columns}[c,onlytextwidth]
    \begin{column}{.46\textwidth}
      {\footnotesize\begin{itemize}\setlength{\itemsep}{.08cm}
        \item Each compute node becomes \(c^{\mathrm{in}}\rightarrow c^{\mathrm{out}}\).
        \item The coordinator becomes the source and the sink.
        \item An edge capacity is how many tokens may cross it per second.
        \item Max flow from source to sink is the serving throughput of this placement.
      \end{itemize}}
      \vspace{.08cm}\takeaway{The graph is an accounting device for one fixed placement.}
    \end{column}
    \begin{column}{.52\textwidth}
      \figureplaceholder{Helix/Helix_Figure_02_Max_Flow_Graph_Abstraction.png}{5.15cm}{Figure 2}
      {\centering\hint{Fig.~2. Capacities are tokens/s.}\par}
    \end{column}
  \end{columns}
  \src{Sec.~4.3, Fig.~2. PDF p.~4.}
\end{frame}

% ----- SLIDE 11 -----
\begin{frame}[t]{Modeling Compute Capacity}
  \smalllabel{Helix / Formulation}\par
  \vspace{.12cm}
  \begin{columns}[T,onlytextwidth]
    \begin{column}{.49\textwidth}
      \stepcard{NODE EDGE}{\(c^{\mathrm{in}}\rightarrow c^{\mathrm{out}}\)}{Capacity is the minimum of profiled compute throughput and the node's network throughput.}
      \vspace{.14cm}
      \stepcard{PROFILING}{Once per node}{Helix measures throughput rather than using the datasheet TFLOP/s from Table 3.}
    \end{column}
    \begin{column}{.49\textwidth}
      \stepcard{LAYER COUNT}{\(T_j\) tokens/s}{\(T_j\) is the profiled rate when that node holds \(j\) layers. It becomes the MILP compute limit.}
      \vspace{.14cm}
      \stepcard{NOT A GPU}{Source and sink}{The coordinator injects tokens and collects tokens. It is not given a compute edge.}
    \end{column}
  \end{columns}
  \vspace{.16cm}\takeaway{Compute capacity is a profiled token rate, and it already takes the node's network limit into account.}
  \src{Sec.~4.3. \(T_j\) is defined with Constraint~3 in Sec.~4.4.}
\end{frame}

% ----- SLIDE 12 -----
\begin{frame}[t]{Modeling Network Capacity}
  \smalllabel{Helix / Formulation}\par
  \vspace{.08cm}
  {\footnotesize A link is valid only in one of three cases:}\par
  \vspace{.08cm}
  \begin{enumerate}\footnotesize
    \setlength{\itemsep}{.06cm}
    \item Coordinator \(\rightarrow\) a node that holds the first layer. Payload is a token (Fig.~2 uses 4 bytes).
    \item A node that holds the last layer \(\rightarrow\) coordinator. Payload is again a token.
    \item Node \(i\rightarrow\) node \(j\) when \(j\) holds the layers needed immediately after \(i\). Payload is an activation (Fig.~2 uses 16 KB).
  \end{enumerate}
  \vspace{.08cm}
  \[
    \text{capacity (tokens/s)}=\frac{\text{profiled average bandwidth}}{\text{bytes per token or activation}}.
  \]
  {\footnotesize The formula is the paper's rule. The byte sizes above are the Fig.~2 annotations, not universal constants.}\par
  \vspace{.06cm}\takeaway{Activations cost far more tokens of budget than the final output token, so wide-area links saturate first.}
  \src{Sec.~4.3. Bandwidth is a one-time average, not a runtime measurement.}
\end{frame}

% ----- SLIDE 13 -----
\begin{frame}[t]{Why Max Flow Represents Serving Throughput}
  \smalllabel{Helix / Formulation}\par
  \vspace{.10cm}
  \begin{itemize}\footnotesize
    \setlength{\itemsep}{.08cm}
    \item One unit of flow is one token that can pass a node or a link in one second.
    \item Conservation forces tokens into \(c^{\mathrm{in}}\) to leave through \(c^{\mathrm{out}}\).
    \item Helix computes the source--sink max flow with preflow-push. That value is the placement's max serving throughput.
    \item The model is a steady token rate. It does not output prompt latency, decode latency, or a percentile.
    \item A later early-stop bound is the sum of node compute rates, divided by the number of layers.
  \end{itemize}
  \vspace{.08cm}\takeaway{Max flow answers ``how many tokens/s can this placement carry?'' It is not a latency model.}
  \src{Sec.~4.3 and the early-stop remark in Sec.~4.5.}
\end{frame}

% ----- SLIDE 14 -----
\begin{frame}[t]{Worked Three-Node Example}
  \smalllabel{Helix / Formulation}\par
  \vspace{.02cm}
  \begin{columns}[c,onlytextwidth]
    \begin{column}{.48\textwidth}
      {\footnotesize\begin{itemize}\setlength{\itemsep}{.05cm}
        \item A100 holds layers 1--2. T4-1 holds layer 1. T4-2 holds layer 3.
        \item Node edges are marked 1500, 1000, and 1000 tokens/s.
        \item Token-sized edges are marked 1.25M, 2.5M, and 625K tokens/s.
        \item Activation edges are marked 458 and 687 tokens/s.
      \end{itemize}}
      \vspace{.06cm}\takeaway{The hundreds, not the millions, set the flow. This is an illustration, not a measured cluster.}
    \end{column}
    \begin{column}{.50\textwidth}
      \figureplaceholder{Helix/Helix_Figure_02_Max_Flow_Graph_Abstraction.png}{5.05cm}{Figure 2}
    \end{column}
  \end{columns}
  \src{Fig.~2 annotations. PDF p.~4. Endpoint of each number is the figure's, not a new measurement.}
\end{frame}

% ----- SLIDE 15 -----
\begin{frame}[t]{MILP Variables and Objective}
  \smalllabel{Helix / MILP}\par
  \vspace{.06cm}
  {\centering\footnotesize
  \begin{tabular}{lll}\toprule
    Symbol & Role & Meaning\\\midrule
    \(s_i\) & integer & first layer on node \(i\)\\
    \(b_{ij}\) & binary & node \(i\) holds exactly \(j\) layers\\
    \(f_{i,j}\) & real & flow from \(i\) to \(j\), in tokens/s\\
    \(d_{i,j}\) & binary & whether link \((i,j)\) is valid\\
    \(\mathit{cond}\) & binary & auxiliaries that linearize overlap\\\bottomrule
  \end{tabular}\par}
  \vspace{.08cm}
  {\footnotesize \(e_i=s_i+\sum_j j\,b_{ij}\). Node \(i\) holds the half-open range \([s_i,e_i)\). \(k\) is the most layers that fit on that GPU.}\par
  \vspace{.06cm}
  \[
    \max \sum_i f_{\mathrm{source},i}
  \]
  \vspace{.02cm}\takeaway{The objective maximizes coordinator outflow. Table 5 groups the variable counts as linear in \(|C|\) or \(|E|\).}
  \src{Sec.~4.4, Tables 5--6. Notation here is compacted from the paper.}
\end{frame}

% ----- SLIDE 16 -----
\begin{frame}[t]{Placement and Memory Constraints}
  \smalllabel{Helix / MILP}\par
  \vspace{.10cm}
  \[
    \sum_{j=1}^{k} b_{ij}=1,
    \qquad 0\le s_i<L,
    \qquad e_i\le L.
  \]
  \begin{itemize}\footnotesize
    \setlength{\itemsep}{.07cm}
    \item Exactly one layer-count is chosen. The range must lie inside the model's \(L\) layers.
    \item Memory enters as \(k\), the maximum number of layers the GPU can hold. Table 6 has no KV-byte variable.
    \item If \(s_j<e_i<e_j\), traffic from \(i\) runs only layers \([e_i,e_j)\) on \(j\). The paper calls this partial inference.
    \item The solver may be run with partial inference disabled. Then a link is valid only when \(e_i=s_j\).
  \end{itemize}
  \vspace{.06cm}\takeaway{GPU memory limits how many layers are placed. KV capacity is enforced later, at schedule time.}
  \src{Sec.~4.4, Constraint~1, and the partial-inference remark.}
\end{frame}

% ----- SLIDE 17 -----
\begin{frame}[t]{Flow and Compute Constraints}
  \smalllabel{Helix / MILP}\par
  \vspace{.12cm}
  {\footnotesize\textbf{Constraint 2: conservation at node \(i\).}}\par
  \[
    \sum_u f_{u,i}=\sum_v f_{i,v}.
  \]
  {\footnotesize\textbf{Constraint 3: profiled inference throughput.}}\par
  \[
    \sum_u f_{u,i}\le \sum_{j=1}^{k} b_{ij}\,T_j.
  \]
  \begin{itemize}\footnotesize
    \setlength{\itemsep}{.06cm}
    \item \(T_j\) is a profiled constant: tokens per second when this node holds \(j\) layers.
    \item The binaries \(b_{ij}\) pick exactly one of those profiled rates.
    \item Flow is still also limited by the network edges on the next slide.
  \end{itemize}
  \vspace{.04cm}\takeaway{A node cannot be assigned more token flow than the profiled rate for the layer count it actually holds.}
  \src{Sec.~4.4, Constraints 2--3. Sums skip node \(i\) itself, as stated in the paper.}
\end{frame}

% ----- SLIDE 18 -----
\begin{frame}[t]{Connectivity and Network Constraints}
  \smalllabel{Helix / MILP}\par
  \vspace{.06cm}
  {\footnotesize\textbf{Constraint 4, compacted.} A coordinator link into \(i\) exists only if \(i\) holds the first layer:}\par
  \[
    s_i\le L\,(1-d_{\mathrm{source},i}).
    \qquad
    L\,d_{i,\mathrm{sink}}\le e_i.
  \]
  {\footnotesize Between compute nodes, \(d_{ij}=1\) only if \(s_j\le e_i<e_j\). Two auxiliary binaries linearize that test, and \(d_{ij}\le 0.5\,\mathit{cond}^{1}+0.5\,\mathit{cond}^{2}\).}\par
  \vspace{.06cm}
  {\footnotesize\textbf{Constraint 5: transmission throughput.}}\par
  \[
    f_{ij}\le d_{ij}\,S_{ij}.
  \]
  {\footnotesize \(S_{ij}\) is the profiled token rate of that link, from Sec.~4.3. Invalid links carry zero flow.}\par
  \vspace{.04cm}\takeaway{Routing variables cannot use a hop that the layer ranges do not justify.}
  \src{Sec.~4.4, Constraints 4--5. The two-auxiliary form is the paper's linearization.}
\end{frame}

% ----- SLIDE 19 -----
\begin{frame}[t]{Making the MILP Practical}
  \smalllabel{Helix / MILP}\par
  \vspace{.10cm}
  \begin{columns}[T,onlytextwidth]
    \begin{column}{.32\textwidth}
      \stepcard{PRUNE}{Average degree 12}{Drop slow links before solving. Sec.~6.8 checks the throughput effect.}
    \end{column}
    \begin{column}{.32\textwidth}
      \stepcard{WARM START}{Heuristic seeds}{Petals, Swarm, and separate pipelines provide the initial solution.}
    \end{column}
    \begin{column}{.32\textwidth}
      \stepcard{EARLY STOP}{Analytic bound}{Stop near the sum of compute rates divided by the layer count.}
    \end{column}
  \end{columns}
  \vspace{.16cm}
  \begin{itemize}\footnotesize
    \setlength{\itemsep}{.06cm}
    \item Relaxing the integers to an LP is rejected: rounding \(s_i\) or \(b_{ij}\) can delete edges and change the flow.
    \item Gurobi stops if 10 minutes pass with no improvement. The budget is 4 hours on a 14-core CPU.
    \item For hundreds of nodes, the paper only suggests cutting the cluster into smaller pieces first.
  \end{itemize}
  \vspace{.06cm}\takeaway{The formulation can certify optimality. The running system also prunes, seeds, and stops early.}
  \src{Sec.~4.5--4.6 and the Helix setup in Sec.~6.2.}
\end{frame}

% ----- SLIDE 20 -----
\begin{frame}[t]{Runtime Components and Control Flow}
  \smalllabel{Helix / Runtime}\par
  \vspace{.12cm}
  \begin{enumerate}\footnotesize
    \setlength{\itemsep}{.08cm}
    \item \textbf{Once per cluster:} profile nodes and links, solve the MILP, load each node's layer range.
    \item \textbf{On request arrival:} the coordinator's scheduler builds one pipeline and sends the request to its first node.
    \item \textbf{On each node:} run the assigned layers for the requests that arrived during the previous batch. Do not add a waiting window.
    \item \textbf{On the last node:} return the output token to the coordinator.
    \item \textbf{On the next token:} reuse that request's pipeline. Do not solve the MILP again.
  \end{enumerate}
  \vspace{.08cm}\takeaway{Placement is a cluster constant. The per-token path is a walk on the graph that placement created.}
  \src{Sec.~5 and Fig.~3. Batching details are in Sec.~5.1.}
\end{frame}

% ----- SLIDE 21 -----
\begin{frame}[t]{Per-Request Pipelines}
  \smalllabel{Helix / Runtime}\par
  \vspace{.10cm}
  \begin{itemize}\footnotesize
    \setlength{\itemsep}{.08cm}
    \item A valid pipeline runs every layer once, in order. Each stage names a node and the layers it will run.
    \item Petals-style fixed pipelines give each pipeline a disjoint set of machines.
    \item Helix lets pipelines intersect. The candidates are the source--sink paths of the placement graph.
    \item After the pipeline is chosen, a node batches every request that arrived while it ran the previous batch. There is no extra collect window.
    \item Later tokens of the same request keep that pipeline.
  \end{itemize}
  \vspace{.08cm}\takeaway{Flexibility comes from many overlapping paths, not from re-solving placement per token.}
  \src{Sec.~5.1. Fixed-pipeline contrast cites prior systems, including HexGen in the related work.}
\end{frame}

% ----- SLIDE 22 -----
\begin{frame}[t]{Converting Max-Flow Weights into Scheduling}
  \smalllabel{Helix / Runtime}\par
  \vspace{.10cm}
  \begin{columns}[T,onlytextwidth]
    \begin{column}{.49\textwidth}
      \stepcard{IWRR}{Interleaved weighted round-robin}{Each vertex has a scheduler. A neighbor's weight is the max-flow value on that outgoing edge.}
      \vspace{.12cm}
      \stepcard{WHY INTERLEAVE}{Avoid bursts}{Frequency follows the weight, but choices are interleaved so one heavy edge is not drained in a burst.}
    \end{column}
    \begin{column}{.49\textwidth}
      \stepcard{THE WALK}{Coordinator, then each node}{Pick stage 1 from the coordinator. Pick stage 2 from that node's scheduler. Stop when the path is complete.}
      \vspace{.12cm}
      \stepcard{MASK}{Before the walk}{Nodes over the KV high-water mark are removed from the candidate list. Details are on the KV slide.}
    \end{column}
  \end{columns}
  \vspace{.14cm}\takeaway{The MILP decides the rates. IWRR is the online rule that realizes those rates.}
  \src{Sec.~5.1. IWRR follows Tabatabaee et al.\ as cited by the paper.}
\end{frame}

% ----- SLIDE 23 -----
\begin{frame}[t]{Pipeline Scheduling Example}
  \smalllabel{Helix / Runtime}\par
  \vspace{.02cm}
  \begin{columns}[c,onlytextwidth]
    \begin{column}{.46\textwidth}
      {\footnotesize\begin{itemize}\setlength{\itemsep}{.07cm}
        \item Vertices are compute nodes. An edge exists only if the MILP placement made it valid.
        \item The number on an edge is flow in the max-flow solution, not link capacity.
        \item The figure's two example pipelines are \(1\rightarrow 5\rightarrow 9\) and \(2\rightarrow 6\rightarrow 8\).
        \item Those paths can share the cluster. They are not two frozen replicas.
      \end{itemize}}
      \vspace{.06cm}\takeaway{Read the integers as scheduled token rates, then read the paths as two requests.}
    \end{column}
    \begin{column}{.52\textwidth}
      \figureplaceholder{Helix/Helix_Figure_04_Per_Request_Pipeline_Scheduling.png}{4.85cm}{Figure 4}
    \end{column}
  \end{columns}
  \src{Sec.~5.1, Fig.~4. PDF p.~8.}
\end{frame}

% ----- SLIDE 24 -----
\begin{frame}[t]{KV-Cache Masking and Backpressure}
  \smalllabel{Helix / Runtime}\par
  \vspace{.12cm}
  \begin{itemize}\footnotesize
    \setlength{\itemsep}{.09cm}
    \item Each GPU has a finite VRAM budget for KV cache. Exceeding it forces offload to host memory, which the paper says significantly cuts throughput.
    \item Output length is unknown until generation ends, so the scheduler cannot know the exact KV footprint.
    \item Helix estimates each node's KV use with the average output length.
    \item A node above the high-water mark is masked out of IWRR until some of its current requests finish.
    \item The paper does not print the numeric threshold or the byte formula for that estimate.
  \end{itemize}
  \vspace{.08cm}\takeaway{Backpressure is a mask on the same IWRR walk. It is not a second placement solver.}
  \src{Sec.~5.2. ``High water mark'' is the paper's term; its value is not specified.}
\end{frame}

% ----- SLIDE 25 -----
\begin{frame}[t]{Implementation and Prototype}
  \smalllabel{Helix / Evaluation}\par
  \vspace{.10cm}
  \begin{columns}[T,onlytextwidth]
    \begin{column}{.49\textwidth}
      \stepcard{PROTOTYPE}{vLLM 0.4.0.post1}{1.5k lines of Python and 1.7k lines of C++. A unified page pool supports partial inference. ZeroMQ carries inter-node traffic. Gurobi solves the MILP.}
      \vspace{.12cm}
      \stepcard{SIMULATOR}{14k lines of Python}{It replays heterogeneous GPUs and links. Geo-distributed and 42-node results use this simulator.}
    \end{column}
    \begin{column}{.49\textwidth}
      \stepcard{FIDELITY}{Single-cluster check}{Against the prototype, decode-throughput error is under 5\%. Prompt-latency error is under 5\%, decode-latency error under 4\%.}
      \vspace{.12cm}
      \stepcard{SCOPE}{Single-GPU nodes}{Multi-GPU nodes are supported by in-node TP, but the prototype experiments use one GPU per node.}
    \end{column}
  \end{columns}
  \vspace{.12cm}\takeaway{Trust the simulator inside the single-cluster error above. Do not treat that as a geo-network error bar.}
  \src{Sec.~6.1 and Sec.~6.3. FP16 execution is stated in Sec.~6.2.}
\end{frame}

% ----- SLIDE 26 -----
\begin{frame}[t]{Evaluation Methodology}
  \smalllabel{Helix / Evaluation}\par
  \vspace{.06cm}
  {\centering\footnotesize
  \begin{tabular}{llll}\toprule
    Setup & Hardware & Network & Where\\\midrule
    Single & 4 A100, 8 L4, 12 T4 & 10 Gb/s, one region & Prototype\\
    Geo & 4 A100; 2 L4+8 T4; 6 L4+4 T4 & \(\sim\)100 Mb/s, \(\sim\)50 ms & Simulator\\
    42-node & 7 node types, including V100 and multi-GPU & same simulator stack & Simulator\\\bottomrule
  \end{tabular}\par}
  \vspace{.08cm}
  \begin{itemize}\footnotesize
    \setlength{\itemsep}{.05cm}
    \item Models: LLaMA-1 30B and LLaMA-2 70B, FP16. Baselines: Swarm, and separate pipelines (SP) inside Helix.
    \item Petals is compared only as a placement (Sec.~6.6), not as an end-to-end server.
    \item Offline metric: average decode throughput, in generated tokens/s. Online adds average prompt and decode latency.
    \item Online arrivals are the Azure shape scaled to 75\% of that cluster's peak. Offline arrivals saturate the cluster.
  \end{itemize}
  \vspace{.04cm}\takeaway{A \(3\times\) on the 42-node simulator is not a \(3\times\) on the 24-node prototype.}
  \src{Sec.~6.2. The 100 Mb/s geo setting is justified by Table 7, whose pairs span 54--204 Mb/s.}
\end{frame}

% ----- SLIDE 27 -----
\begin{frame}[t]{Workload Characteristics}
  \smalllabel{Helix / Evaluation}\par
  \vspace{.02cm}
  \begin{columns}[c,onlytextwidth]
    \begin{column}{.48\textwidth}
      {\footnotesize\begin{itemize}\setlength{\itemsep}{.07cm}
        \item Source: Azure Conversation. Requests with input \(>2048\) or output \(>1024\) are removed.
        \item What remains: 16{,}657 requests, mean input 763 tokens, mean output 232 tokens.
        \item Online: 30 s warmup, then 30 min. Offline: 1 min warmup, then 10 min.
        \item Fig.~5 is the trace's length mix and arrival shape, before the 75\% scaling used online.
      \end{itemize}}
      \vspace{.06cm}\takeaway{Latency ratios below are for this pruned trace, not for short-prompt serving.}
    \end{column}
    \begin{column}{.50\textwidth}
      \figureplaceholder{Helix/Helix_Figure_05_Azure_Conversation_Statistics.png}{4.55cm}{Figure 5}
    \end{column}
  \end{columns}
  \src{Sec.~6.2, Fig.~5. PDF p.~8.}
\end{frame}

% ----- SLIDE 28 -----
\begin{frame}[t]{Single-Cluster Results}
  \smalllabel{Helix / Results}\par
  \vspace{.02cm}
  \figureplaceholder{Helix/Helix_Figure_06_Single_Cluster_Results.png}{2.58cm}{Figure 6}
  {\centering\hint{Fig.~6. Prototype. (a--d) decode throughput; (e--h) online latency. Box 25--75, whisker 5--95, red line = median.}\par}
  \vspace{.02cm}
  {\centering\scriptsize
  \begin{tabular}{llcccc}\toprule
    Model & Versus & Offline tput & Online tput & Prompt lat. & Decode lat.\\\midrule
    30B & SP & \(+4\%\) & \(+14\%\) & \(-2\%\) & \(+10\%\)\\
    30B & Swarm & \(2.14\times\) & \(2.07\times\) & \(-32\%\) & \(-12\%\)\\
    70B & SP & \(1.86\times\) & \(1.69\times\) & \multicolumn{2}{c}{average latency higher}\\
    70B & Swarm & \(1.94\times\) & \(2.00\times\) & \(-15\%\) & \(+16\%\)\\\bottomrule
  \end{tabular}\par}
  \vspace{.02cm}
  {\scriptsize 30B: three replicas, one GPU type each. 70B vs.\ SP: latency is higher because SP serves most requests on A100s and leaves L4/T4 idle.}\par
  \vspace{.02cm}\takeaway{Throughput gains are real on the prototype. They are not a uniform latency win.}
  \src{Sec.~6.3, Fig.~6. Pairs are offline then online. PDF p.~9.}
\end{frame}

% ----- SLIDE 29 -----
\begin{frame}[t]{Geo-Distributed Results}
  \smalllabel{Helix / Results}\par
  \vspace{.02cm}
  \figureplaceholder{Helix/Helix_Figure_07_Geo_Distributed_Cluster_Results.png}{2.55cm}{Figure 7}
  {\centering\hint{Fig.~7. Simulator. H = Helix, S = Swarm. Same GPUs as the single cluster, with \(\sim\)100 Mb/s between clusters.}\par}
  \vspace{.02cm}
  {\footnotesize\textbf{LLaMA 30B vs.\ SP:} \(+7\%\) / \(+10\%\) throughput; prompt \(-14\%\), decode latency \(+2\%\). \textbf{Vs.\ Swarm:} \(2.41\times\) and \(2.33\times\) across the two modes; prompt \(-66\%\), decode latency \(-24\%\).}\par
  \vspace{.02cm}
  {\footnotesize\textbf{LLaMA 70B.} Helix's pipeline is 28\% shallower than Swarm, and 19\% shallower than Helix's own fast-network plan. Vs.\ SP: \(1.61\times\) / \(1.79\times\). Vs.\ Swarm: \(1.92\times\) / \(1.97\times\); prompt \(-21\%\), decode \(-7\%\). Swarm's offline prompt latency is 71 s, \(7.5\times\) Helix.}\par
  \vspace{.02cm}\takeaway{The abstract's ``up to 66\% and 24\%'' is this 30B online comparison against Swarm, in simulation.}
  \src{Sec.~6.4, Fig.~7. Sec.~6.4 does not label which of 2.41 and 2.33 is online. PDF p.~10.}
\end{frame}

% ----- SLIDE 30 -----
\begin{frame}[t]{High GPU-Heterogeneity Results}
  \smalllabel{Helix / Results}\par
  \vspace{.02cm}
  \begin{columns}[c,onlytextwidth]
    \begin{column}{.48\textwidth}
      {\footnotesize\begin{itemize}\setlength{\itemsep}{.05cm}
        \item Simulator, LLaMA 70B, 42 nodes, 7 types. V100, T4, and \(2\times\)T4 cannot form a pipeline alone.
        \item SP drops those machines. SP+ adds a mixed pipeline that uses them.
        \item Vs.\ Swarm / SP / SP+, offline throughput is \(1.37\times\), \(2.91\times\), \(2.24\times\).
        \item Online: \(1.48\times\), \(3.29\times\), \(2.54\times\). Prompt latency changes by \(-17\%\), \(-1\%\), \(-7\%\).
        \item Helix's decode latency is slightly higher: baselines leave slow GPUs idle.
      \end{itemize}}
      \vspace{.04cm}\takeaway{The abstract's \(3.3\times\) is the \(3.29\times\) online gain versus SP, rounded.}
    \end{column}
    \begin{column}{.50\textwidth}
      \figureplaceholder{Helix/Helix_Figure_08_High_GPU_Heterogeneity_Results.png}{4.70cm}{Figure 8}
    \end{column}
  \end{columns}
  \src{Sec.~6.5, Fig.~8. PDF p.~10. Not a prototype measurement.}
\end{frame}

% ----- SLIDE 31 -----
\begin{frame}[t]{Model Placement Deep Dive}
  \smalllabel{Helix / Ablation}\par
  \vspace{.02cm}
  \figureplaceholder{Helix/Helix_Figure_09_Model_Placement_Comparison.png}{2.85cm}{Figure 9}
  {\centering\hint{Fig.~9. Offline LLaMA 70B. Every method uses Helix's scheduler. Integers in (b) are layer counts.}\par}
  \vspace{.03cm}
  \begin{itemize}\footnotesize
    \setlength{\itemsep}{.04cm}
    \item Vs.\ Petals / Swarm: \(1.23\times\) / \(2.10\times\) on the single cluster, and \(1.49\times\) / \(2.38\times\) on the geo-distributed simulator.
    \item Single-cluster case: Swarm's tail is four T4s with four layers each, so A100 and L4 stay under-used.
    \item Petals leaves eight T4s and one L4 under-used. Helix's plan keeps almost every node busy.
  \end{itemize}
  \vspace{.03cm}\takeaway{The placement gap remains after the scheduler is held fixed.}
  \src{Sec.~6.6, Fig.~9. PDF p.~11. Petals is not an end-to-end baseline.}
\end{frame}

% ----- SLIDE 32 -----
\begin{frame}[t]{Request Scheduling Deep Dive}
  \smalllabel{Helix / Ablation}\par
  \vspace{.02cm}
  \figureplaceholder{Helix/Helix_Figure_10_Request_Scheduling_Comparison.png}{2.72cm}{Figure 10}
  {\centering\hint{Fig.~10. Offline LLaMA 70B. Every method uses Helix's placement. (b) is the geo-distributed case study.}\par}
  \vspace{.03cm}
  \begin{itemize}\footnotesize
    \setlength{\itemsep}{.04cm}
    \item Vs.\ Swarm / random: \(+30\%\) / \(+29\%\) on the single cluster, and \(+22\%\) / \(+15\%\) on the geo-distributed clusters.
    \item Vs.\ shortest queue, geo-distributed only: \(+19\%\) decode throughput.
    \item In (b), prompt traffic waits 5--16 s on three slow links. One of those queues is caused by a scheduler three hops away.
  \end{itemize}
  \vspace{.03cm}\takeaway{A local queue rule cannot see a congested cut that starts several hops upstream.}
  \src{Sec.~6.7, Fig.~10. PDF p.~12.}
\end{frame}

% ----- SLIDE 33 -----
\begin{frame}[t]{MILP Optimization Ablation}
  \smalllabel{Helix / Ablation}\par
  \vspace{.02cm}
  \begin{columns}[c,onlytextwidth]
    \begin{column}{.46\textwidth}
      {\footnotesize\begin{itemize}\setlength{\itemsep}{.05cm}
        \item Offline LLaMA 70B. 24-node is the geo setup; 42-node is the high-heterogeneity setup.
        \item Degree-12 pruning drops 50\% and 72\% of the links, and 36\% and 46\% of the problem size.
        \item With pruning, decode throughput is 16\% and 2\% higher. The paper attributes this to an easier search under a time limit.
        \item Warm start finds the same placement, in 43\% and 8\% less wall time.
      \end{itemize}}
    \end{column}
    \begin{column}{.52\textwidth}
      \figureplaceholder{Helix/Helix_Figure_11_MILP_Optimization_Ablation.png}{3.15cm}{Figure 11}
      {\centering\scriptsize
      \begin{tabular}{lcc}\toprule
        & With pruning & Without\\\midrule
        24-node & 876 var, 1122 cstr & 1376 var, 1848 cstr\\
        42-node & 2144 var, 2772 cstr & 4004 var, 5502 cstr\\\bottomrule
      \end{tabular}\par}
    \end{column}
  \end{columns}
  \vspace{.04cm}\takeaway{Pruning changes the search problem. It is not a proof that the deleted links were useless.}
  \src{Sec.~6.8, Fig.~11, Table 8. PDF p.~12. Speedups vary by instance, as the paper notes.}
\end{frame}

% ----- SLIDE 34 -----
\begin{frame}[t]{Solver Time and Solution Quality}
  \smalllabel{Helix / Ablation}\par
  \vspace{.02cm}
  \begin{columns}[c,onlytextwidth]
    \begin{column}{.48\textwidth}
      {\footnotesize\begin{itemize}\setlength{\itemsep}{.07cm}
        \item Instance: LLaMA 30B on 4 L4s and 6 T4s. This is smaller than the 24- and 42-node runs.
        \item The best feasible placement reaches the optimum in under 5 minutes.
        \item Certifying that optimum takes more than an hour of branch-and-bound.
        \item The red dotted line is that cluster's optimal throughput. The upper bound tightens as cuts are added.
      \end{itemize}}
      \vspace{.06cm}\takeaway{Early stop is reasonable here. Fig.~12 does not certificate the 42-node plan.}
    \end{column}
    \begin{column}{.50\textwidth}
      \figureplaceholder{Helix/Helix_Figure_12_MILP_Solution_Quality.png}{4.85cm}{Figure 12}
    \end{column}
  \end{columns}
  \src{Sec.~6.9, Fig.~12. PDF p.~13.}
\end{frame}

% ----- SLIDE 35 -----
\begin{frame}[t]{Helix vs.\ ThunderServe, Hetis, and Tessera}
  \smalllabel{Cross-paper comparison}\par
  \vspace{.04cm}
  {\centering\scriptsize
  \renewcommand{\arraystretch}{1.12}
  \begin{tabularx}{\textwidth}{@{}>{\raggedright\arraybackslash}p{1.55cm}*{4}{>{\raggedright\arraybackslash}X}@{}}\toprule
    & ThunderServe & Hetis & Helix & Tessera\\\midrule
    Granularity & Prefill / decode phase & Dense module vs.\ attention head & Layer range and pipeline & CUDA kernel\\
    Decision & Group, phase, route & Primary set, then heads & Placement, then a path & Kernel-to-GPU map\\
    Time scale & Offline search + light reschedule & Offline plan + online re-dispatch & Once per cluster + per request & Offline MILP + monitor window\\
    State moved & KV at the phase boundary & QKV, attention output, KV shards & Tokens and activations & Cut tensors; KV delta\\
    Adaptation & Phase role and failure & Re-dispatch, migration & KV mask; path is reused & Latency / throughput switch\\
    Environment & Heterogeneous cloud & Local heterogeneous cluster & Heterogeneous, including geo & Heterogeneous pairs and clusters\\
    Models & LLM & LLM & LLM (LLaMA 30B/70B) & LLM, SSM, MLLM, diffusion\\
    Main limit & Heuristic phase search & Head groups; local network & One-shot plan; geo is simulated & MILP ignores GPU memory\\\bottomrule
  \end{tabularx}\par}
  \vspace{.06cm}\takeaway{These systems cut different bottlenecks. Their speedups are not a ranking.}
  \src{Our synthesis. Helix cells: this PDF. Other cells: each system's own paper, not Helix's related work.}
\end{frame}

% ----- SLIDE 36 -----
\begin{frame}[t]{Limitations and Threats to Validity}
  \smalllabel{Helix / Discussion}\par
  \vspace{.08cm}
  \begin{itemize}\footnotesize
    \setlength{\itemsep}{.07cm}
    \item Placement is computed once. A workload shift does not trigger a new MILP in the described runtime.
    \item Geo-distributed and 42-node numbers are simulated. The \(<5\%\) error is a single-cluster prototype check.
    \item The KV mask uses average output length. The paper does not specify the high-water mark.
    \item Pruning and the 10-minute stall rule mean a deployed plan need not be a proven optimum.
    \item Decode latency can rise when Helix actually uses the slow GPUs that SP leaves idle.
    \item Evidence is FP16 LLaMA-1 30B and LLaMA-2 70B on one pruned Azure trace. The conclusion's ``guaranteed optimal'' is the MILP's claim, not every timed run.
  \end{itemize}
  \vspace{.06cm}\takeaway{Report the cluster, the simulator, and the metric together with any ratio.}
  \src{Sec.~4.5, Sec.~6, and Sec.~8. The last bullet contrasts the conclusion with Sec.~6.2 and Sec.~6.9.}
\end{frame}

% ----- SLIDE 37 -----
\begin{frame}[t]{Conclusion}
  \smalllabel{Helix / Conclusion}\par
  \vspace{.14cm}
  \begin{columns}[T,onlytextwidth]
    \begin{column}{.32\textwidth}
      \stepcard{1}{A placement is a flow}{Edge capacity is tokens per second. Max flow is the throughput of that placement, not its latency.}
    \end{column}
    \begin{column}{.32\textwidth}
      \stepcard{2}{Two clocks}{MILP places layers once. IWRR turns the flow into a per-request path and reuses it for later tokens.}
    \end{column}
    \begin{column}{.32\textwidth}
      \stepcard{3}{Where it pays}{The large gains appear when one GPU type cannot hold the model, or when the wide-area network is the bottleneck.}
    \end{column}
  \end{columns}
  \vspace{.20cm}
  \begin{itemize}\footnotesize
    \setlength{\itemsep}{.06cm}
    \item On the reported settings, decode throughput improves by up to \(3.29\times\) versus SP (abstract: \(3.3\times\)).
    \item The largest latency cuts versus Swarm, 66\% prompt and 24\% decode, are the simulated geo-distributed LLaMA 30B online run.
  \end{itemize}
  \vspace{.08cm}\takeaway{Helix jointly chooses layer placement and request paths. It does not claim a latency win on every baseline.}
  \src{Sec.~8, with the ratios taken from Sec.~6.4 and Sec.~6.5 rather than from the abstract alone.}
\end{frame}

% ----- SLIDE 38 -----
\begin{frame}
  \vfill
  \begin{center}
    {\fontsize{32}{38}\selectfont\bfseries\color{Ink}Questions \& Discussion}\par
    \vspace{.34cm}
    {\large\color{Accent}Max flow \quad / \quad per-request pipelines \quad / \quad what is simulated}\par
    \vspace{.42cm}
    {\small\color{Muted}Helix (layer and path) \quad next to phase, head, and kernel systems}
  \end{center}
  \vfill
\end{frame}

\end{document}



